We study the \emph{Radical Identity Testing} problem (RIT): Given an algebraic circuit representing a polynomial $f\in \mathbb{Z}[x_1, \dots, x_k]$ and nonnegative integers $a_1, \dots, a_k$ and $d_1, \dots,$ $d_k$, written in binary, test whether  the polynomial vanishes at the \emph{real radicals} $\mysqrt{d_1}{a_1}, \dots,\mysqrt{d_k}{a_k}$, i.e., test whether $f(\mysqrt{d_1}{a_1}, \dots, \mysqrt{d_k}{a_k}) = 0$. We place the problem in  {\coNP} assuming the Generalised Riemann Hypothesis (GRH), improving on the straightforward {\PSPACE} upper bound obtained by reduction to the existential theory of reals. Next we consider a restricted version, called $2$-RIT, where the radicals are square roots of prime numbers, written in binary. It was known since the work of Chen and Kao~\cite{chen-kao} that $2$-RIT is at least as hard as the polynomial identity testing problem, however no better upper bound than {\PSPACE} was known prior to our work. We show that $2$-RIT is in {\coRP} assuming GRH and in {\coNP} unconditionally. 
Our proof relies on theorems from algebraic and analytic number theory, such as the Chebotarev density theorem and quadratic reciprocity.





